% !TeX root = ../初等数学笔记.tex
% 由主文档通过 \input 载入；请编译主文档。
\setcounter{chapter}{4} % 保留原书章号。
\chapter{方程与方程组}
\label{chap:em-equations}
方程求解的目标是确定全部解，而不是把式子变得好看。每一次操作都应回答两个问题：原解是否仍被保留，新式是否产生原来不允许的解。以下除特别说明外，未知数均在实数范围内取值。

\section{方程、解集与等价变形}
\emph{方程}（equation）是要求未知数满足的等式，解集由定义域内使等式成立的全部值组成。两方程解集相同才称等价。在公共定义域内，同加同减保持等价；同乘一个可能为零的式子只保证原解仍成立，同除该式子则可能丢解。

平方通常是单向操作：$u=v$ 推出 $u^2=v^2$，逆向只能得到 $u=v$ 或 $u=-v$。若已经知道 $u,v\ge0$，平方才等价。求解时可以使用必要条件得到有限个候选值，但必须代回原方程筛选；验根是恢复充分性的步骤，并非装饰性检查。

\section{一元一次方程与线性方程组}
对实参数 $a,b$，方程 $ax=b$ 在 $a\ne0$ 时唯一解为 $b/a$，在 $a=0,b\ne0$ 时无解，在 $a=b=0$ 时每个实数都是解。参数退化必须在除法前处理。

方程组的解要同时满足各方程。交换方程、将某一方程乘以非零常数、将一方程的倍数加到另一方程上，都是可逆操作，因而保持公共解集。对
\[
 ax+by=e,\qquad cx+dy=f,
\]
消元给出 $(ad-bc)x=ed-bf$。若 $ad-bc\ne0$，再回代可得
\[
 x=\frac{ed-bf}{ad-bc},\qquad y=\frac{af-ec}{ad-bc}.
\]
若 $ad-bc=0$，应继续检查方程是矛盾还是相关，而不能使用上述商式。一般线性方程组也按逐个消元、识别矛盾或自由变量、最后回代的顺序处理。

\section{一元二次方程及根与系数}
设 $a\ne0$。将 $ax^2+bx+c=0$ 乘以 $4a$ 并配方，得到 $(2ax+b)^2=b^2-4ac$。记\emph{判别式}（discriminant）为 $\Delta=b^2-4ac$，则
\begin{equation}\label{eq:em-quadratic}
 x=\frac{-b\pm\sqrt\Delta}{2a}\quad(\Delta\ge0).
\end{equation}
$\Delta>0$ 有两个不同实根，$\Delta=0$ 有一个二重根，$\Delta<0$ 没有实根。若两根按重数记为 $x_1,x_2$，展开 $a(x-x_1)(x-x_2)$ 并比较系数得
\[
 x_1+x_2=-\frac ba,\qquad x_1x_2=\frac ca.
\]
这就是二次方程的\emph{韦达关系}（Viète's formulas）。反之，和为 $s$、积为 $p$ 的两实数必须是 $t^2-st+p=0$ 的根，所以存在的条件是 $s^2-4p\ge0$。

\section{可化归的高次、分式与根式方程}
高次方程先观察可分解或换元的结构。例如 $x^4-5x^2+4=0$ 令 $t=x^2\ge0$，得 $(t-1)(t-4)=0$，故 $x=\pm1,\pm2$。新变量的范围与每个新值对应的原变量都不能省略。

分式方程先排除分母零点，在剩余定义域内乘以公分母才等价。例如 $(x^2-1)/(x-1)=0$ 的候选根为 $\pm1$，但 $1$ 不在定义域，故只有 $-1$。根式方程 $\sqrt{x+2}=x$ 要求 $x\ge0$，平方得 $(x-2)(x+1)=0$，只有 $x=2$ 合法。若在一开始没有记录右侧非负，就容易留下增根 $-1$。

\section{含参数方程与非线性方程组}
分类依据应是运算或解的性质发生变化的位置。例如 $(a-1)x=a^2-1$ 在 $a=1$ 时恒成立，在 $a\ne1$ 时唯一解为 $a+1$。若提前约去 $a-1$，就丢失了整个退化情形。

根的位置可以结合判别式与韦达关系判断。实系数二次方程的两实根均为正，当且仅当判别式非负、根和为正、根积为正；后两条件说明两根同号且不可能同时为负。非线性方程组也可利用对称关系：$x+y=5,xy=6$ 意味着 $x,y$ 是 $t^2-5t+6=0$ 的两根，因此解为 $(2,3),(3,2)$。这种方法先确定无序根，再恢复有序数对。

\section{例题、习题与学习检查}
\begin{example}
求使 $x^2-2mx+m+2=0$ 有两个不同正实根的实数 $m$。
\end{example}
\begin{solve}
不同实根要求 $\Delta=4(m-2)(m+1)>0$，即 $m<-1$ 或 $m>2$；根和为 $2m>0$，根积为 $m+2>0$。三条件取交集得 $m>2$。$m=2$ 虽给出正根 $2$，但为重根，不符合题设。
\end{solve}
\begin{enumerate}
\item 解 $\sqrt{2x+3}=x$。答：$x=3$。
\item 解 $x+y=7,\ x^2+y^2=25$。答：$(3,4),(4,3)$，先由平方和求积。
\item 判断 $x(x-1)=x$ 中两侧约去 $x$ 会漏掉什么。答：原解为 $0,2$，约去后漏掉 $0$。
\end{enumerate}

